%%%%%%%%%%%%%%%%%%%%%%%%%%%%%%%%%%%%%%%%%%%%%%%%%%%%%%%%%%%%%%%%%%%%%%%%%%%%%%%
%% Capítulo 2 --- Fundamentação teórica
%%
%% Capítulo expositivo: apresenta os conceitos que os capítulos seguintes
%% pressupõem, na ordem da cadeia de medição (fenômeno físico -> aquisição ->
%% avaliação normativa -> decisão -> implantação). O confronto com a
%% literatura fica no capítulo de trabalhos relacionados.
%%
%% Cada seção termina aplicando o conceito ao dispositivo estudado, com
%% números que o leitor pode refazer. Os valores de projeto vêm de:
%%   rolamento e rotação   docs/relatorio/relatorio-kaelix.tex, tab:parametros
%%   enlace LoRa           src/comms/lora.cpp (SF9, 125 kHz, 4/5, preâmbulo 8)
%%   carga útil            src/power/sleep.cpp (21 bytes)
%%   rigidez e massa       Capítulo de resultados, Seção do invólucro
%%   MPU6050               folha de dados rev. 3.1 e mapa de registradores
%%                         rev. 4.0, conferidos no documento primário
%%%%%%%%%%%%%%%%%%%%%%%%%%%%%%%%%%%%%%%%%%%%%%%%%%%%%%%%%%%%%%%%%%%%%%%%%%%%%%%

\makeatletter
\@ifundefined{toprule}{%
  \newcommand{\toprule}{\hline}%
  \newcommand{\midrule}{\hline}%
  \newcommand{\bottomrule}{\hline}%
}{}
\makeatother

\chapter{Fundamentação teórica}
\label{cap:fundamentacao}

Este capítulo apresenta os conceitos sobre os quais o restante do trabalho se apoia. A
ordem de exposição acompanha o percurso da informação no dispositivo estudado: parte do
fenômeno físico --- a vibração de uma máquina rotativa e as assinaturas que cada tipo de
falha deixa nela ---, passa pela conversão desse fenômeno em sequência de números, pelos
critérios com que a norma avalia a severidade da vibração e pelos métodos com que um
algoritmo decide se a máquina se afastou da operação normal, e termina nas restrições de
um equipamento embarcado: comunicação sem fio, energia e temperatura. Ao final de cada
seção, o conceito é aplicado ao dispositivo estudado, de modo que os números usados nos
capítulos seguintes possam ser refeitos a partir daqui.

%%%%%%%%%%%%%%%%%%%%%%%%%%%%%%%%%%%%%%%%%%%%%%%%%%%%%%%%%%%%%%%%%%%%%%%%%%%%%%%
\section{Manutenção de máquinas rotativas}
\label{sec:fund-manutencao}

A manutenção de equipamentos industriais é usualmente classificada em três estratégias,
distinguidas pelo momento em que a intervenção ocorre \cite{mobley2002}. Na manutenção
\emph{corretiva}, o equipamento opera até falhar e só então é reparado; o custo de
planejamento é nulo, mas a falha ocorre em momento não escolhido e pode danificar
componentes que estavam íntegros. Na manutenção \emph{preventiva}, a intervenção ocorre em
intervalos fixos de tempo ou de uso, definidos a partir da vida média esperada dos
componentes; a falha inesperada torna-se menos frequente, mas componentes em bom estado são
substituídos sem necessidade, e falhas que ocorrem antes do intervalo continuam não
detectadas. Na manutenção \emph{preditiva}, ou baseada em condição, a intervenção é decidida
pelo estado medido do equipamento: acompanha-se uma grandeza física que reflete sua
condição e intervém-se quando ela indica degradação.

\citeonline{jardine2006} descrevem um programa de manutenção baseada em condição em três
etapas: aquisição de dados, processamento dos dados e tomada de decisão de manutenção. A
divisão é útil porque separa problemas de natureza distinta. A aquisição é um problema de
instrumentação: que grandeza medir, com que sensor, em que ponto da máquina e com que taxa.
O processamento é um problema de tratamento de sinais: como extrair, dos dados brutos,
indicadores que respondam à condição do equipamento. A decisão é um problema de
classificação ou de detecção: a partir desses indicadores, decidir se a condição exige
intervenção. Os capítulos seguintes adotam uma divisão análoga, em que a qualidade de cada
etapa limita a das seguintes.

Entre as grandezas empregadas no monitoramento de máquinas rotativas, a vibração é a mais
difundida, porque a maior parte dos defeitos mecânicos --- desbalanceamento, desalinhamento,
folga, desgaste de rolamentos e de engrenagens --- altera as forças dinâmicas no interior da
máquina, e essas forças se transmitem à carcaça, onde podem ser medidas sem interromper a
operação \cite{randall2011livro}.

%%%%%%%%%%%%%%%%%%%%%%%%%%%%%%%%%%%%%%%%%%%%%%%%%%%%%%%%%%%%%%%%%%%%%%%%%%%%%%%
\section{Fundamentos de vibração mecânica}
\label{sec:fund-vibracao}

\subsection{Deslocamento, velocidade e aceleração}

Vibração é o movimento oscilatório de um ponto em torno de uma posição de equilíbrio. O
movimento pode ser descrito por três grandezas relacionadas: o deslocamento $x(t)$, a
velocidade $v(t) = \mathrm{d}x/\mathrm{d}t$ e a aceleração $a(t) = \mathrm{d}v/\mathrm{d}t$.
Para um movimento senoidal de frequência $f$ e amplitude de deslocamento $X$, com
$\omega = 2\pi f$,

\begin{equation}
x(t) = X \sin(\omega t), \qquad
v(t) = \omega X \cos(\omega t), \qquad
a(t) = -\omega^2 X \sin(\omega t).
\label{eq:fund-senoide}
\end{equation}

As amplitudes das três grandezas diferem, portanto, por fatores de $\omega$: para um mesmo
deslocamento, a velocidade cresce proporcionalmente à frequência, e a aceleração, ao
quadrado dela. Essa relação tem duas consequências práticas. A primeira é que a grandeza
escolhida para avaliação pondera as frequências de maneira diferente: o deslocamento
realça as componentes de baixa frequência, a aceleração realça as de alta frequência, e a
velocidade ocupa posição intermediária. A segunda é que, conhecida uma das grandezas, as
outras podem ser obtidas por derivação ou integração; em particular, um acelerômetro, que
mede $a(t)$, pode fornecer a velocidade por integração, operação discutida na
Seção~\ref{sec:fund-integracao}.

\subsection{Descritores escalares de um sinal}
\label{sec:fund-descritores}

Um sinal de vibração registrado durante um intervalo é frequentemente resumido por
descritores escalares. Para uma janela de $N$ amostras $x_n$, com média $\bar{x}$ e
desvio padrão $\sigma$, os mais usados são \cite{randall2011livro}:

\begin{equation}
x_{\mathrm{rms}} = \sqrt{\frac{1}{N}\sum_{n=1}^{N} x_n^2},
\qquad
x_{\mathrm{pico}} = \max_n |x_n|,
\qquad
\mathrm{FC} = \frac{x_{\mathrm{pico}}}{x_{\mathrm{rms}}},
\qquad
\kappa = \frac{\frac{1}{N}\sum_{n=1}^{N} (x_n-\bar{x})^4}{\sigma^4}.
\label{eq:fund-descritores}
\end{equation}

O valor eficaz ($x_{\mathrm{rms}}$) mede a energia média do sinal e é a base da avaliação
normativa de severidade (Seção~\ref{sec:fund-norma}). O valor de pico e o pico a pico
(diferença entre o maior e o menor valor da janela) medem a excursão máxima. O fator de
crista (FC) e a curtose ($\kappa$) medem a \emph{impulsividade}: quanto o sinal concentra
sua energia em eventos breves e intensos. Valores de referência ajudam a interpretá-los.
Uma senoide pura tem $\mathrm{FC} = \sqrt{2} \approx \num{1,41}$ e $\kappa = \num{1,5}$;
um ruído gaussiano tem $\kappa = 3$. Um sinal composto por impactos espaçados, como o
produzido por um defeito localizado em rolamento (Seção~\ref{sec:fund-rolamento}), tem
fator de crista e curtose bem maiores, porque poucos valores extremos dominam a média de
quarta potência. Esses descritores são insensíveis, porém, à \emph{frequência} em que a
energia se concentra e à relação entre componentes de frequências diferentes; essa
informação só é acessível pela análise espectral (Seção~\ref{sec:fund-espectro}).

\subsection{Sistema de um grau de liberdade e ressonância}
\label{sec:fund-1gdl}

O modelo mais simples de uma estrutura que vibra é o sistema de um grau de liberdade: uma
massa $m$ ligada a uma base por uma mola de rigidez $k$ e um amortecedor. Deslocada do
equilíbrio e liberada, a massa oscila na frequência natural \cite{brandt2011}

\begin{equation}
f_n = \frac{1}{2\pi}\sqrt{\frac{k}{m}}.
\label{eq:fund-fn}
\end{equation}

Quando a base vibra com frequência $f$, a razão entre a amplitude do movimento da massa e a
da base é a \emph{transmissibilidade}. Com a razão de frequências $r = f/f_n$ e a razão de
amortecimento $\zeta$,

\begin{equation}
T(r) = \sqrt{\frac{1 + (2\zeta r)^2}{(1-r^2)^2 + (2\zeta r)^2}}.
\label{eq:fund-transmissibilidade}
\end{equation}

Para $r \ll 1$, $T \approx 1$: a massa acompanha a base. À medida que $f$ se aproxima de
$f_n$, a transmissibilidade cresce, e em $r = 1$ é limitada apenas pelo amortecimento ---
é a ressonância. Acima de $r = \sqrt{2}$, $T < 1$ e o movimento da base é atenuado.

Esse modelo descreve um problema central da medição de vibração: o sensor não está preso
diretamente à superfície da máquina, mas a um conjunto de peças --- base, fixação,
invólucro, placa --- que tem rigidez finita. Esse conjunto e a massa que ele sustenta
formam um sistema com frequência natural própria, a \emph{frequência de montagem}, e o
sensor mede o movimento da massa, não o da máquina. Para que a medição reproduza a
vibração da superfície, a transmissibilidade precisa ficar próxima de 1 em toda a banda de
interesse. Com amortecimento desprezível, o erro de amplitude é $1/(1-r^2) - 1$, e o
critério de erro máximo de \SI{10}{\percent} exige

\begin{equation}
\frac{1}{1-r^2} \le \num{1,1}
\quad\Longrightarrow\quad
r \le \frac{1}{\sqrt{11}} \approx \num{0,30}
\quad\Longrightarrow\quad
f_n \ge \num{3,3}\, f_{\max},
\label{eq:fund-criterio-montagem}
\end{equation}

\noindent
em que $f_{\max}$ é a maior frequência que se deseja medir. Para a banda normativa, cujo
limite superior é \SI{1000}{\hertz} (Seção~\ref{sec:fund-norma}), a frequência de
montagem precisa estar acima de cerca de \SI{3,3}{\kilo\hertz}.

\paragraph{Aplicação ao dispositivo estudado.} Com a rigidez do caminho de montagem de
\SI{1,45e6}{\newton\per\meter} e a massa de \SI{118}{\gram} obtidas na análise do
invólucro, a Equação~\eqref{eq:fund-fn} dá $f_n \approx \SI{557}{\hertz}$. O erro de
amplitude atinge \SI{10}{\percent} em $r = \num{0,30}$, isto é, em torno de
\SI{168}{\hertz}, e em \SI{500}{\hertz} ($r \approx \num{0,90}$) a transmissibilidade sem
amortecimento é de cerca de 5.

%%%%%%%%%%%%%%%%%%%%%%%%%%%%%%%%%%%%%%%%%%%%%%%%%%%%%%%%%%%%%%%%%%%%%%%%%%%%%%%
\section{Assinaturas de falha em máquinas rotativas}
\label{sec:fund-falhas}

Cada tipo de defeito altera a vibração de um modo característico, e é essa assinatura que
permite, a partir do sinal medido, inferir o que está errado na máquina. As assinaturas
dividem-se em dois grupos, conforme a escala de frequência em que se manifestam.

\subsection{Defeitos síncronos com a rotação}

Desbalanceamento e desalinhamento produzem forças que se repetem a cada volta do eixo, e
por isso se manifestam em múltiplos inteiros da frequência de rotação $f_r$, as
\emph{harmônicas} da rotação \cite{randall2011livro}. No desbalanceamento, o centro de
massa do rotor não coincide com o eixo de rotação, e a força centrífuga resultante gira com
o eixo; a assinatura dominante é a componente em $1\times f_r$, cuja amplitude cresce com a
massa desbalanceada. No desalinhamento, os eixos acoplados não são colineares, e o
acoplamento impõe esforços que variam mais de uma vez por volta; a assinatura típica é o
aumento da componente em $2\times f_r$ em relação à de $1\times f_r$.

Como as rotações industriais usuais estão entre algumas centenas e alguns milhares de
rotações por minuto, essas componentes situam-se entre poucos hertz e algumas centenas de
hertz. Um motor a \SI{1750}{rpm}, por exemplo, tem $f_r = \SI{29,17}{\hertz}$, e suas três
primeiras harmônicas estão abaixo de \SI{100}{\hertz}.

\subsection{Defeitos em rolamentos}
\label{sec:fund-rolamento}

O rolamento de esferas é composto por uma pista interna, solidária ao eixo, uma pista
externa, fixa à carcaça, e um conjunto de esferas mantidas equidistantes por uma gaiola.
Um defeito localizado --- uma trinca ou um lascamento em uma das pistas --- é atingido
periodicamente pelas esferas, e cada passagem produz um impacto. A taxa de repetição desses
impactos depende apenas da geometria do rolamento e da rotação. Para um defeito na pista
externa e na pista interna, respectivamente \cite{randall2011},

\begin{equation}
\mathrm{BPFO} = \frac{N_b}{2}\, f_r \left(1 - \frac{d}{D_p}\cos\theta\right),
\qquad
\mathrm{BPFI} = \frac{N_b}{2}\, f_r \left(1 + \frac{d}{D_p}\cos\theta\right),
\label{eq:fund-bpfo}
\end{equation}

\noindent
com $N_b$ o número de esferas, $d$ o diâmetro da esfera, $D_p$ o diâmetro primitivo e
$\theta$ o ângulo de contato. Para o rolamento SKF~6205 ($N_b = 9$, $d/D_p = \num{0,203}$,
$\theta = 0$) a \SI{1750}{rpm}, as expressões dão $\mathrm{BPFO} = \SI{104,6}{\hertz}$ e
$\mathrm{BPFI} = \SI{157,9}{\hertz}$.

A frequência de repetição, contudo, não é onde a energia do defeito aparece. Cada impacto é
um evento breve, que excita as ressonâncias da estrutura entre o rolamento e o ponto de
medição, tipicamente na faixa de alguns quilohertz. O sinal resultante é uma sequência de
oscilações amortecidas na frequência de ressonância, que se repetem na taxa BPFO ou BPFI.
Em termos espectrais, a energia do defeito concentra-se em torno da ressonância, e a
frequência característica aparece como \emph{modulação} dessa portadora de alta
frequência \cite{randall2011}. Nos estágios iniciais do defeito, as linhas espectrais na
própria BPFO tendem a ser fracas diante das demais componentes da máquina; a informação
diagnóstica está na alta frequência.

É por isso que o procedimento consolidado de diagnóstico de rolamentos é a \emph{análise de
envelope}. O sinal é filtrado em uma banda em torno da ressonância; calcula-se sua
envoltória, isto é, a amplitude instantânea, pelo módulo do sinal analítico obtido com a
transformada de Hilbert; e o espectro dessa envoltória revela linhas na frequência de
repetição dos impactos \cite{randall2011}. A escolha da banda de filtragem é a etapa
crítica, e a curtose espectral --- a curtose calculada em função da frequência --- é a
ferramenta estabelecida para localizar a banda em que o sinal é mais impulsivo
\cite{antoni2006,antoni2007}. Todo esse procedimento pressupõe que a banda de ressonância
tenha sido adquirida: um sistema de medição cuja banda termina abaixo da ressonância não
dispõe do sinal sobre o qual a análise de envelope opera.

%%%%%%%%%%%%%%%%%%%%%%%%%%%%%%%%%%%%%%%%%%%%%%%%%%%%%%%%%%%%%%%%%%%%%%%%%%%%%%%
\section{Aquisição e processamento digital de sinais}
\label{sec:fund-sinais}

\subsection{Amostragem, frequência de Nyquist e rebatimento}
\label{sec:fund-amostragem}

Um sistema digital não registra o sinal contínuo $a(t)$, mas amostras
$a_n = a(n/f_s)$ tomadas à taxa de amostragem $f_s$. O teorema da amostragem estabelece
que um sinal pode ser reconstruído a partir de suas amostras se não contiver componentes
acima de $f_s/2$, a \emph{frequência de Nyquist} \cite{brandt2011}. Uma componente acima
desse limite não desaparece na amostragem: suas amostras são idênticas às de uma
componente de frequência mais baixa, dentro da faixa $[0, f_s/2]$. O fenômeno chama-se
\emph{rebatimento} (\emph{aliasing}), e a frequência aparente de uma componente em $f$ é

\begin{equation}
f_{\mathrm{ap}} = \left| f - f_s \cdot \mathrm{round}\!\left(\frac{f}{f_s}\right) \right|.
\label{eq:fund-alias}
\end{equation}

Depois da amostragem, a componente rebatida é indistinguível de uma componente legítima
naquela frequência, e nenhum processamento posterior consegue separá-las. A única proteção
é remover o conteúdo acima de $f_s/2$ \emph{antes} da amostragem, por um filtro
passa-baixas chamado filtro \emph{anti-rebatimento} (\emph{anti-aliasing}). Pela mesma
razão, reduzir a taxa de um sinal já digitalizado --- a \emph{decimação} --- exige filtrar
o sinal abaixo da nova frequência de Nyquist antes de descartar amostras \cite{brandt2011}.

\paragraph{Aplicação ao dispositivo estudado.} Com $f_s = \SI{1}{\kilo\hertz}$, a
frequência de Nyquist é \SI{500}{\hertz}. Uma ressonância de rolamento em
\SI{3,5}{\kilo\hertz} que chegasse sem filtro ao conversor seria observada, pela
Equação~\eqref{eq:fund-alias}, em $|3500 - 1000 \cdot 4| = \SI{500}{\hertz}$, uma
frequência em que nada existe fisicamente.

\subsection{Análise espectral}
\label{sec:fund-espectro}

A representação de um sinal em função da frequência é obtida pela transformada discreta
de Fourier (DFT), calculada na prática pelo algoritmo da transformada rápida (FFT). Para
uma janela de $N$ amostras,

\begin{equation}
X_k = \sum_{n=0}^{N-1} x_n\, e^{-\mathrm{j}2\pi k n / N}, \qquad k = 0, 1, \dots, N-1,
\label{eq:fund-dft}
\end{equation}

\noindent
e o coeficiente $X_k$ corresponde à frequência $k\,f_s/N$ \cite{brandt2011}. Duas
propriedades limitam o que o espectro mostra. A \emph{resolução} em frequência é
$\Delta f = f_s/N$, o inverso da duração da janela: duas componentes mais próximas que
$\Delta f$ não são separadas. E a janela finita provoca \emph{vazamento espectral}: uma
componente cuja frequência não coincide com um múltiplo de $\Delta f$ espalha energia
pelos coeficientes vizinhos, efeito que se atenua multiplicando o sinal por uma função de
janela que o leve suavemente a zero nas bordas, como a de Hann \cite{brandt2011}. Para
sinais aleatórios, o espectro de uma única janela é uma estimativa muito ruidosa, e
usa-se a \emph{densidade espectral de potência}, obtida pela média dos espectros de várias
janelas.

Para o dispositivo estudado, com $f_s = \SI{1}{\kilo\hertz}$ e janelas de \num{512}
amostras, a janela dura \SI{0,512}{\second} e a resolução é de cerca de
\SI{1,95}{\hertz}, suficiente para separar as harmônicas de uma rotação de
\SI{29}{\hertz}.

\subsection{Integração de aceleração para velocidade}
\label{sec:fund-integracao}

Como a norma avalia a severidade em velocidade (Seção~\ref{sec:fund-norma}) e o
acelerômetro fornece aceleração, é preciso integrar o sinal. A integração pode ser feita de
duas formas \cite{brandt2014}.

No \emph{domínio do tempo}, a velocidade é aproximada pela soma acumulada das amostras de
aceleração multiplicadas pelo intervalo de amostragem. O método é simples, mas tem um
defeito estrutural: qualquer componente contínua presente na aceleração --- um desvio de
zero do sensor, por exemplo --- é integrada como uma rampa, e o erro cresce linearmente com
a duração do registro. Um filtro passa-alta aplicado antes da integração reduz esse efeito,
mas introduz seus próprios transitórios nas bordas da janela.

No \emph{domínio da frequência}, usa-se a relação da Equação~\eqref{eq:fund-senoide}: cada
componente de frequência $f$ da velocidade é a componente correspondente da aceleração
dividida por $\mathrm{j}2\pi f$. Calcula-se a DFT da aceleração, divide-se cada coeficiente
por $\mathrm{j}2\pi f_k$, anulam-se os coeficientes fora da banda de interesse --- o que
elimina a componente contínua e a divisão por zero em $f = 0$ --- e retorna-se ao tempo pela
transformada inversa. \citeonline{brandt2014} compararam os dois métodos e concluíram que a
integração no domínio da frequência é a mais exata para registros longos.

Como exemplo, uma aceleração senoidal de amplitude \SI{3,0}{\meter\per\second\squared} a
\SI{50}{\hertz} corresponde a uma velocidade de amplitude
$3{,}0/(2\pi \cdot 50) \approx \SI{9,55}{\milli\meter\per\second}$ e valor eficaz
$\num{9,55}/\sqrt{2} \approx \SI{6,75}{\milli\meter\per\second}$. É esse valor analítico que
serve de referência para comparar os métodos de integração no
Capítulo~\ref{cap:resultados}.

%%%%%%%%%%%%%%%%%%%%%%%%%%%%%%%%%%%%%%%%%%%%%%%%%%%%%%%%%%%%%%%%%%%%%%%%%%%%%%%
\section{Transdutores de vibração}
\label{sec:fund-transdutores}

\subsection{Acelerômetros piezoelétricos e MEMS}

O transdutor de referência no monitoramento de máquinas é o acelerômetro piezoelétrico:
uma massa sísmica apoiada sobre um cristal que gera carga elétrica proporcional à força, e
portanto à aceleração. Sua frequência de ressonância interna é elevada, o que lhe confere
resposta plana até vários quilohertz, e seu ruído é baixo; em contrapartida, exige
condicionamento de sinal dedicado e tem custo alto \cite{albarbar2008}.

Os acelerômetros de sistemas microeletromecânicos (MEMS) capacitivos são fabricados pelos
mesmos processos dos circuitos integrados. Uma microestrutura suspensa por molas
desloca-se sob aceleração, e o deslocamento altera a capacitância entre eletrodos. O
conversor analógico-digital e o processamento digital ficam no mesmo encapsulamento, e o
sensor entrega amostras digitais por barramento serial. O custo é de uma a duas ordens de
grandeza menor que o de um piezoelétrico, mas a banda e o ruído dependem fortemente do
modelo: há MEMS projetados para monitoramento de condição, com banda de dezenas de
quilohertz, e MEMS de consumo, projetados para orientação e movimento humano, com banda de
algumas centenas de hertz \cite{albarbar2008,fidali2024}. Os critérios que decidem se um
MEMS serve para uma aplicação de vibração são a banda útil, a densidade espectral de ruído,
a faixa de medição e a estabilidade do zero.

A densidade espectral de ruído, expressa em $\si{\micro g}/\sqrt{\si{\hertz}}$, indica o
ruído por unidade de banda. Para ruído branco, o ruído eficaz na banda $B$ é a densidade
multiplicada por $\sqrt{B}$: quanto mais larga a banda analisada, maior o ruído total, e
menor a menor vibração que o sensor consegue distinguir.

\subsection{O acelerômetro MPU6050}
\label{sec:fund-mpu6050}

O acelerômetro adotado no dispositivo estudado é o MPU6050, uma unidade de medição inercial
de consumo que combina acelerômetro e giroscópio de três eixos. As características
relevantes para medição de vibração, conforme a folha de dados e o mapa de registradores do
fabricante, estão na Tabela~\ref{tab:fund-mpu6050}.

\begin{table}[!htb]
\centering
\caption{Características do acelerômetro do MPU6050 relevantes para medição de vibração
\cite{invensense2011ps,invensense2012rm}.}
\label{tab:fund-mpu6050}
\small
\begin{tabular}{@{}ll@{}}
\toprule
Característica & Valor \\
\midrule
Faixa de medição                         & $\pm2$, $\pm4$, $\pm8$ ou $\pm16$~g, programável \\
Resolução do conversor                   & \num{16} bits \\
Taxa de saída de dados                   & \SIrange{4}{1000}{\hertz}, programável \\
Filtro passa-baixas digital (DLPF)       & \SIrange{5}{260}{\hertz}, programável \\
Densidade espectral de ruído             & \SI{400}{\micro g}/$\sqrt{\si{\hertz}}$ a \SI{10}{\hertz} \\
Tolerância de zero (eixos X e Y)         & $\pm\SI{50}{\milli g}$ \\
\bottomrule
\end{tabular}
\end{table}

Dois pontos da tabela condicionam o uso do sensor em vibração. O primeiro é a banda. A taxa
máxima de saída do acelerômetro é de \SI{1}{\kilo\hertz}, o que situa a frequência de
Nyquist em \SI{500}{\hertz}; mas a banda efetiva é fixada pelo filtro passa-baixas digital
interno, configurado pelo campo \texttt{DLPF\_CFG}, e não pela taxa de saída. Os valores
possíveis estão na Tabela~\ref{tab:fund-dlpf}. Na configuração de banda mais larga,
\texttt{DLPF\_CFG}~=~0, que é também o valor após a inicialização, a banda do acelerômetro
é de \SI{260}{\hertz}; não há configuração em que o acelerômetro opere sem esse filtro
\cite{invensense2012rm}. O mapa de registradores informa a banda e o atraso de cada
configuração, mas não a ordem do filtro nem sua atenuação acima da banda, de modo que o
grau de proteção contra rebatimento não é especificado pelo fabricante.
%% ATENÇÃO: os capítulos de desenvolvimento e resultados tratam a banda útil
%% como 500 Hz (Nyquist) e modelam uma cadeia "sem filtro" como o estado do
%% sensor quando o DLPF não é configurado. Pela Tabela tab:fund-dlpf, o
%% acelerômetro nunca opera sem DLPF, e a banda máxima é 260 Hz. Revisar esses
%% capítulos antes da versão final.

\begin{table}[!htb]
\centering
\caption{Banda e atraso do filtro passa-baixas digital do acelerômetro do MPU6050, por
configuração do campo \texttt{DLPF\_CFG}, com taxa de saída de \SI{1}{\kilo\hertz}
\cite{invensense2012rm}.}
\label{tab:fund-dlpf}
\small
\begin{tabular}{@{}ccc@{}}
\toprule
\texttt{DLPF\_CFG} & Banda (\si{\hertz}) & Atraso (\si{\milli\second}) \\
\midrule
0 & 260 & \num{0}    \\
1 & 184 & \num{2,0}  \\
2 & 94  & \num{3,0}  \\
3 & 44  & \num{4,9}  \\
4 & 21  & \num{8,5}  \\
5 & 10  & \num{13,8} \\
6 & 5   & \num{19,0} \\
\bottomrule
\end{tabular}
\end{table}

O segundo ponto é o ruído. Tratando a densidade de \SI{400}{\micro g}/$\sqrt{\si{\hertz}}$
como ruído branco, o ruído eficaz na banda de \SI{260}{\hertz} é de
$400 \cdot \sqrt{260} \approx \SI{6,5}{\milli g}$, cerca de
\SI{0,063}{\meter\per\second\squared}. A tolerância de zero de $\pm\SI{50}{\milli g}$
significa, por sua vez, que o sensor pode apresentar uma componente contínua de até
aproximadamente \SI{0,5}{\meter\per\second\squared} sem estar fora de especificação ---
exatamente o tipo de desvio que a integração no domínio do tempo acumula
(Seção~\ref{sec:fund-integracao}).

%%%%%%%%%%%%%%%%%%%%%%%%%%%%%%%%%%%%%%%%%%%%%%%%%%%%%%%%%%%%%%%%%%%%%%%%%%%%%%%
\section{Avaliação normativa da severidade de vibração}
\label{sec:fund-norma}

A avaliação da vibração de máquinas industriais medida em partes não rotativas é
normalizada pela ISO 20816-3:2022, que substituiu a ISO 10816-3:2009
\cite{iso208163_2022,iso108163_2009}. A norma aplica-se a máquinas com potência acima de
\SI{15}{\kilo\watt} e rotação entre \num{120} e \num{30000}~r/min, e adota como grandeza de
avaliação a \emph{velocidade eficaz de banda larga}. A justificativa está no próprio texto
normativo: vibrações com o mesmo valor eficaz de velocidade na faixa de \SIrange{10}{1000}{\hertz}
podem, em geral, ser consideradas de mesma severidade, ao passo que critérios em
deslocamento ou em aceleração variariam com a frequência \cite{iso208163_2022}. Pela mesma
razão, a norma exige do instrumento resposta plana em pelo menos
\SIrange{10}{1000}{\hertz}, com limite inferior de até \SI{2}{\hertz} para máquinas com
rotação próxima ou abaixo de \num{600}~r/min, e determina que o sinal de acelerômetros
montados na carcaça seja integrado para fornecer velocidade. A norma de instrumentação
correspondente torna verificável o requisito de resposta plana, com tolerância de cerca de
\SI{10}{\percent} na banda passante \cite{iso2954_2012}.

O valor medido é comparado com quatro zonas de avaliação:

\begin{description}
  \item[Zona A.] Vibração típica de máquinas recém-comissionadas.
  \item[Zona B.] Vibração aceitável para operação de longo prazo sem restrição.
  \item[Zona C.] Vibração insatisfatória para operação contínua de longo prazo; a máquina
        pode operar por período limitado, até que haja oportunidade de ação corretiva.
  \item[Zona D.] Vibração suficientemente severa para causar dano à máquina.
\end{description}

Os limites entre as zonas dependem do porte da máquina e do tipo de fundação, conforme a
Tabela~\ref{tab:fund-zonas}. A fundação é dita rígida quando a primeira frequência natural
do conjunto máquina--fundação está suficientemente acima da frequência de excitação
principal, e flexível no caso contrário. Os valores de zona são idênticos nas edições de
2009 e de 2022.

\begin{table}[!htb]
\centering
\caption{Limites entre zonas de avaliação, em velocidade eficaz de banda larga
(\si{\milli\meter\per\second}), para máquinas industriais \cite{iso208163_2022}. Grupo 1:
potência acima de \SI{300}{\kilo\watt}, ou máquinas elétricas com altura de eixo de pelo
menos \SI{315}{\milli\meter}. Grupo 2: potência de \SIrange{15}{300}{\kilo\watt}, ou
máquinas elétricas com altura de eixo entre \num{160} e \SI{315}{\milli\meter}.}
\label{tab:fund-zonas}
\small
\begin{tabular}{@{}llccc@{}}
\toprule
Grupo & Fundação & A/B & B/C & C/D \\
\midrule
1 & rígida   & \num{2,3} & \num{4,5} & \num{7,1}  \\
1 & flexível & \num{3,5} & \num{7,1} & \num{11,0} \\
2 & rígida   & \num{1,4} & \num{2,8} & \num{4,5}  \\
2 & flexível & \num{2,3} & \num{4,5} & \num{7,1}  \\
\bottomrule
\end{tabular}
\end{table}

A tabela torna concreta a sensibilidade da avaliação ao erro de medição. A velocidade
eficaz de \SI{6,75}{\milli\meter\per\second} do exemplo da Seção~\ref{sec:fund-integracao}
está na zona B para máquinas do grupo 1 em fundação flexível, na zona C para o grupo 1
rígido e o grupo 2 flexível, e na zona D para o grupo 2 rígido. Um erro de medição de
poucos milímetros por segundo, portanto, muda a classificação e a recomendação de
manutenção dela decorrente.

%%%%%%%%%%%%%%%%%%%%%%%%%%%%%%%%%%%%%%%%%%%%%%%%%%%%%%%%%%%%%%%%%%%%%%%%%%%%%%%
\section{Detecção de anomalias}
\label{sec:fund-anomalias}

\subsection{Paradigmas de aprendizado}

A decisão sobre a condição da máquina pode ser formulada de modos diferentes, conforme os
dados disponíveis para treinar o modelo \cite{chandola2009}. Na abordagem
\emph{supervisionada}, o modelo aprende a partir de exemplos rotulados de cada condição ---
normal e cada tipo de falha --- e classifica novos dados em uma dessas classes. Na
abordagem \emph{semissupervisionada}, ou de classe única, o modelo aprende apenas a partir
de exemplos da condição normal e marca como anomalia o que se afasta dela. Na abordagem
\emph{não supervisionada}, não há rótulo algum, e o modelo supõe que as anomalias são raras
e diferentes da maioria dos dados.

Em manutenção preditiva, a abordagem supervisionada esbarra em um problema prático: uma
máquina recém-instalada não tem histórico de falhas, e reunir exemplos rotulados de cada
defeito exige que eles ocorram. O treino apenas com dados de operação normal, disponíveis
desde o primeiro dia, é por isso o modo de operação natural de um dispositivo instalado
em uma máquina sem histórico \cite{chandola2009}. O preço é que o modelo detecta
\emph{afastamento} da normalidade, e não identifica \emph{qual} defeito está presente.

\subsection{Isolation Forest}
\label{sec:fund-iforest}

O Isolation Forest \cite{liu2008} parte de uma observação simples: anomalias são poucas e
diferentes, e por isso são mais fáceis de \emph{isolar} do resto dos dados do que os pontos
normais. O algoritmo constrói um conjunto de árvores de isolamento. Cada árvore é
construída sobre uma subamostra aleatória dos dados de treino, dividindo-a recursivamente:
em cada nó, escolhe-se ao acaso um atributo e um valor de corte entre o mínimo e o máximo
desse atributo, até que cada ponto fique isolado em uma folha ou que se atinja um limite de
profundidade. Um ponto anômalo, por estar em região de baixa densidade, tende a ser isolado
em poucas divisões, isto é, perto da raiz; um ponto normal exige mais divisões.

O comprimento do caminho $h(x)$ entre a raiz e a folha em que o ponto $x$ cai mede,
portanto, quão fácil é isolá-lo. Para tornar essa medida comparável entre árvores
construídas com subamostras de tamanho $\psi$, ela é normalizada pelo comprimento médio
esperado do caminho em uma árvore binária de busca com $\psi$ pontos,

\begin{equation}
c(\psi) = 2H(\psi-1) - \frac{2(\psi-1)}{\psi}, \qquad H(i) \approx \ln i + \num{0,5772},
\label{eq:fund-cn}
\end{equation}

\noindent
e o escore de anomalia é

\begin{equation}
s(x, \psi) = 2^{-\,\mathbb{E}[h(x)] / c(\psi)},
\label{eq:fund-score}
\end{equation}

\noindent
em que $\mathbb{E}[h(x)]$ é a média do comprimento do caminho sobre todas as árvores
\cite{liu2008}. O escore está entre 0 e 1. Valores próximos de 1 indicam ponto facilmente
isolável, candidato a anomalia; valores bem abaixo de \num{0,5} indicam ponto normal; e, se
todos os pontos têm escore próximo de \num{0,5}, os dados não contêm anomalias distintas. Os
valores recomendados pelos autores são $t = 100$ árvores e subamostras de $\psi = 256$
pontos, para as quais $c(256) \approx \num{10,2}$.

O algoritmo tem características que o tornam atraente para execução embarcada: o treino
opera sobre subamostras pequenas, a inferência percorre cada árvore uma única vez com uma
comparação por nó, e o modelo treinado é inteiramente descrito por atributos e valores de
corte \cite{liu2008}. A viabilidade de treiná-lo e executá-lo em microcontrolador foi
demonstrada \cite{antonini2023}.

\subsection{Limiar de decisão e taxa de falso alarme}
\label{sec:fund-limiar}

O escore contínuo de anomalia precisa ser convertido em decisão binária por um
\emph{limiar}: acima dele, o ponto é marcado como anômalo. O limiar não é propriedade do
modelo, e sim uma escolha que troca um tipo de erro pelo outro. Um limiar baixo detecta mais
falhas e marca mais pontos normais como anomalia; um limiar alto faz o contrário.

Na implementação de referência do algoritmo, a escolha é controlada pelo parâmetro de
contaminação, que representa a fração esperada de anomalias nos dados de treino. No valor
padrão, o limiar corresponde à referência interpretativa de \num{0,5} do escore; com um valor
numérico, o limiar é posicionado de modo que essa fração dos pontos de treino fique acima
dele. Em ambos os casos, o parâmetro altera apenas o deslocamento do limiar, e não o escore
calculado para cada ponto \cite{scikitlearn2026}.

Quando o treino contém apenas operação normal, o limiar pode ser calibrado diretamente pela
taxa de falso alarme desejada: se o limiar é o quantil de ordem $1-\alpha$ dos escores de
dados normais não usados no ajuste das árvores, espera-se que uma fração $\alpha$ da
operação normal fique acima dele. Uma alternativa com fundamentação estatística mais forte
modela apenas a cauda da distribuição dos escores pela teoria de valores extremos e fixa o
limiar a partir de um risco alvo \cite{siffer2017}.

%%%%%%%%%%%%%%%%%%%%%%%%%%%%%%%%%%%%%%%%%%%%%%%%%%%%%%%%%%%%%%%%%%%%%%%%%%%%%%%
\section{Avaliação de detectores}
\label{sec:fund-avaliacao}

\subsection{Matriz de confusão e taxas derivadas}

Com a anomalia como classe positiva, cada decisão do detector cai em uma de quatro
categorias: verdadeiro positivo (VP), anomalia corretamente detectada; falso positivo (FP),
operação normal marcada como anomalia; verdadeiro negativo (VN); e falso negativo (FN),
anomalia não detectada. Dessas contagens derivam as taxas usuais \cite{fawcett2006}:

\begin{equation}
\mathrm{TPR} = \frac{\mathrm{VP}}{\mathrm{VP} + \mathrm{FN}},
\qquad
\mathrm{FPR} = \frac{\mathrm{FP}}{\mathrm{FP} + \mathrm{VN}},
\qquad
\mathrm{precisão} = \frac{\mathrm{VP}}{\mathrm{VP} + \mathrm{FP}},
\qquad
F_1 = \frac{2\cdot \mathrm{precisão} \cdot \mathrm{TPR}}{\mathrm{precisão} + \mathrm{TPR}}.
\label{eq:fund-taxas}
\end{equation}

A taxa de verdadeiros positivos (TPR) é a \emph{taxa de detecção}; a taxa de falsos
positivos (FPR) é a \emph{taxa de falso alarme}. Em monitoramento contínuo, a segunda tem
peso operacional particular: um dispositivo que avalia a máquina a cada poucos minutos e
marca \SI{10}{\percent} da operação normal como anomalia gera dezenas de alarmes falsos por
dia, e alarmes frequentes demais deixam de ser atendidos. A acurácia --- fração de decisões
corretas --- é pouco informativa quando as classes são desbalanceadas, porque um detector
que nunca dispara alcança acurácia alta se as anomalias são raras.

\subsection{Curva ROC, área sob a curva e área parcial}

Como o limiar é uma escolha, um detector não tem uma única TPR, mas uma para cada limiar. A
curva ROC traça a TPR contra a FPR à medida que o limiar varre todos os valores possíveis,
e a área sob essa curva (AUC) resume o desempenho em um número entre \num{0,5}, o de uma
decisão ao acaso, e 1, o de um detector perfeito \cite{fawcett2006}. A AUC tem uma
interpretação direta: é a probabilidade de que uma anomalia escolhida ao acaso receba escore
maior que um ponto normal escolhido ao acaso. Por depender apenas da ordenação dos escores,
ela é invariante a qualquer deslocamento do limiar.

Essa mesma propriedade é sua limitação para um dispositivo de campo, que opera em um único
ponto da curva, em falso alarme baixo. Dois detectores com a mesma AUC podem ter
desempenhos muito diferentes nessa região. A \emph{área parcial sob a curva} (pAUC)
restringe a integração à faixa de FPR de interesse, entre 0 e um limite $\beta$
\cite{mcclish1989}. Na forma normalizada adotada no desafio de referência de detecção não
supervisionada de anomalias em máquinas \cite{koizumi2020},

\begin{equation}
\mathrm{pAUC}(\beta) = \frac{1}{\beta}\int_{0}^{\beta} \mathrm{TPR}(\mathrm{FPR})\; \mathrm{d}\,\mathrm{FPR},
\label{eq:fund-pauc}
\end{equation}

\noindent
em que a divisão por $\beta$ normaliza o resultado para o intervalo $[0, 1]$; o desafio
adota $\beta = \num{0,1}$.

\subsection{Validação cruzada e vazamento de informação}
\label{sec:fund-validacao}

O desempenho de um modelo só é informativo se medido em dados não usados no treino. Na
\emph{validação cruzada em $k$ partições}, os dados são divididos em $k$ partes; o modelo é
treinado $k$ vezes, cada vez com uma parte reservada para teste, e o desempenho é a média
das $k$ avaliações, acompanhada de sua dispersão. A dispersão entre partições é informação,
e não ruído a descartar: ela indica quanto o resultado depende de quais dados foram usados
no teste.

A validade da validação cruzada exige que as partes de treino e de teste sejam
independentes. Quando não são, ocorre \emph{vazamento de informação}: o modelo tem acesso,
no treino, a informação que não estaria disponível na aplicação real, e o desempenho
medido superestima o desempenho verdadeiro \cite{kaufman2012}. Em sinais de vibração, a
forma mais comum de vazamento vem da segmentação. Um ensaio de alguns segundos é dividido em
muitas janelas, e janelas do mesmo ensaio compartilham características que nada têm a ver
com a falha --- nível de ruído, ganho do sensor, condição de fixação. Se as janelas são
distribuídas aleatoriamente entre treino e teste, o modelo pode acertar reconhecendo o
ensaio, e não a falha. \citeonline{varejao2025} denominam esse efeito \emph{viés de
similaridade}. A correção é a \emph{validação por grupos}: todas as janelas de um mesmo
ensaio são mantidas na mesma partição, de modo que o teste é sempre feito sobre ensaios
inteiramente desconhecidos pelo modelo.

%%%%%%%%%%%%%%%%%%%%%%%%%%%%%%%%%%%%%%%%%%%%%%%%%%%%%%%%%%%%%%%%%%%%%%%%%%%%%%%
\section{Sistemas embarcados para monitoramento}
\label{sec:fund-embarcados}

Um nó sensor de monitoramento é um sistema embarcado: um computador dedicado a uma única
função, integrado ao equipamento que atende e sujeito a restrições de memória, de
processamento e, sobretudo, de energia. O componente central é o microcontrolador, que
reúne processador, memória e periféricos --- conversores, interfaces seriais, temporizadores
--- em um único circuito integrado.

O microcontrolador adotado no dispositivo estudado é o ESP32-S3, com dois núcleos de
processamento de até \SI{240}{\mega\hertz}, \SI{512}{\kilo\byte} de memória SRAM,
coprocessadores de ultrabaixo consumo e rádios Wi-Fi e Bluetooth integrados
\cite{espressif2026}. Para a operação a bateria, a característica decisiva são os modos de
baixo consumo: no modo \emph{deep-sleep}, os núcleos principais e a maior parte dos
periféricos são desligados, e o consumo cai para cerca de \SI{7}{\micro\ampere}
\cite{espressif2026}.

A estratégia usual de um nó a bateria é o \emph{ciclo de trabalho}: o dispositivo permanece
a maior parte do tempo em modo de baixo consumo, desperta em intervalos regulares, adquire
e processa uma janela de sinal, transmite o resultado e volta a dormir. A corrente média
é a média das correntes de cada fase, ponderada pela fração do ciclo que ela ocupa, e a
autonomia é a capacidade da bateria dividida por essa corrente média. Como a fase de sono
ocupa quase todo o ciclo, a corrente de sono e a duração das fases ativas --- em especial a
transmissão --- dominam a autonomia.

Executar a detecção no próprio dispositivo, em vez de transmitir o sinal bruto para
processamento remoto, é o que viabiliza esse ciclo: transmitir uma janela de \num{512}
amostras de três eixos custaria ordens de grandeza mais tempo de rádio que transmitir um
escore. A execução de modelos de aprendizado de máquina em microcontroladores, com memória
de centenas de quilobytes e consumo de miliwatts, é o campo denominado TinyML
\cite{antonini2023,barbariol2022}.

%%%%%%%%%%%%%%%%%%%%%%%%%%%%%%%%%%%%%%%%%%%%%%%%%%%%%%%%%%%%%%%%%%%%%%%%%%%%%%%
\section{Comunicação de longo alcance: LoRa}
\label{sec:fund-lora}

LoRa é uma técnica de modulação de rádio projetada para enlaces de longo alcance e baixo
consumo, à custa de baixa taxa de dados \cite{augustin2016}. Ela usa espalhamento
espectral por \emph{chirp}: cada símbolo é representado por um sinal cuja frequência varia
linearmente ao longo de toda a largura de banda do canal. O espalhamento permite
demodular sinais abaixo do nível de ruído, o que aumenta o alcance. LoRa designa a camada
física; LoRaWAN é um protocolo de rede definido sobre ela, com servidores de rede e
gerência de dispositivos. Um enlace ponto a ponto, entre o dispositivo e um receptor
próprio, pode usar a modulação LoRa sem o protocolo LoRaWAN.

Três parâmetros controlam o compromisso entre alcance, taxa e consumo \cite{semtech2015}:

\begin{description}
  \item[Fator de espalhamento (SF),] de 6 a 12: cada símbolo carrega SF bits e é formado
        por $2^{\mathrm{SF}}$ \emph{chips}. Cada incremento de SF dobra a duração do
        símbolo e aumenta a sensibilidade do receptor.
  \item[Largura de banda (BW),] tipicamente \SI{125}{\kilo\hertz}: banda mais larga
        encurta o símbolo e reduz a sensibilidade.
  \item[Taxa de codificação (CR),] de 4/5 a 4/8: fração de redundância acrescentada para
        correção de erros.
\end{description}

A taxa de símbolos é $R_s = \mathrm{BW}/2^{\mathrm{SF}}$, e a duração do símbolo é
$T_s = 1/R_s$. A duração total de um pacote, o \emph{tempo no ar}, é a soma da duração do
preâmbulo e da carga útil \cite{semtech2015}:

\begin{equation}
\begin{split}
T_{\mathrm{pacote}} = {} & (n_{\mathrm{pre}} + \num{4,25})\,T_s \\
& + \left[ 8 + \max\!\left( \left\lceil \frac{8\,\mathrm{PL} - 4\,\mathrm{SF} + 28 + 16\,\mathrm{CRC} - 20\,\mathrm{IH}}{4(\mathrm{SF} - 2\,\mathrm{DE})} \right\rceil (\mathrm{CR} + 4),\; 0 \right) \right] T_s,
\end{split}
\label{eq:fund-toa}
\end{equation}

\noindent
em que $n_{\mathrm{pre}}$ é o número de símbolos do preâmbulo, PL o número de bytes da carga
útil, CRC vale 1 quando a verificação de redundância está habilitada, IH vale 1 quando o
cabeçalho é omitido, DE vale 1 quando a otimização para baixa taxa está ativa, e CR vale de
1 (para 4/5) a 4 (para 4/8).

\paragraph{Aplicação ao dispositivo estudado.} O enlace do dispositivo usa SF9,
BW~=~\SI{125}{\kilo\hertz}, CR~=~4/5, preâmbulo de 8 símbolos e carga útil de \num{21}
bytes, com cabeçalho e CRC habilitados e sem otimização para baixa taxa. A duração do
símbolo é $2^9/\num{125000} = \SI{4,096}{\milli\second}$; o preâmbulo ocupa
$(8 + \num{4,25}) \times \num{4,096} \approx \SI{50,2}{\milli\second}$; a carga útil ocupa
$8 + \lceil 176/36 \rceil \times 5 = 33$ símbolos, ou \SI{135,2}{\milli\second}. O tempo
no ar resultante é de cerca de \SI{185}{\milli\second}. Em SF12, a mesma carga útil
ocuparia o rádio por um tempo várias vezes maior, e é por isso que a escolha do fator de
espalhamento é, ao mesmo tempo, uma escolha de alcance e de autonomia.

%%%%%%%%%%%%%%%%%%%%%%%%%%%%%%%%%%%%%%%%%%%%%%%%%%%%%%%%%%%%%%%%%%%%%%%%%%%%%%%
\section{Transferência de calor e limites térmicos}
\label{sec:fund-termica}

Um dispositivo fixado à carcaça de um motor recebe calor da superfície em que está montado.
Em regime permanente, a condução de calor por uma camada de espessura $L$, área $A$ e
condutividade térmica $k_t$ pode ser descrita por uma \emph{resistência térmica}
$R = L/(k_t A)$, e a diferença de temperatura através dela é o produto do fluxo de calor
pela resistência, $\Delta T = q\,R$ \cite{bergman2011}. Camadas em série somam suas
resistências, como resistores elétricos. O modelo permite estimar a temperatura interna do
dispositivo a partir da temperatura da carcaça, das resistências do caminho e da potência
dissipada internamente.

A consequência para a escolha de material é direta: um invólucro de polímero, cuja
condutividade é duas a três ordens de grandeza menor que a de um metal, oferece resistência
térmica muito maior e isola o interior da carcaça quente; um invólucro metálico faz o
contrário. Como a Seção~\ref{sec:fund-1gdl} mostrou, a rigidez --- e portanto a
frequência de montagem --- também depende do material, e os dois requisitos podem apontar
em sentidos opostos.

Entre os componentes internos, a bateria é tipicamente o mais sensível à temperatura. A
temperatura acelera os mecanismos de degradação das células de íon-lítio, e a carga em
temperatura elevada é particularmente prejudicial, razão pela qual os fabricantes
especificam faixas de temperatura mais estreitas para carga do que para descarga
\cite{ma2018}. O limite térmico de um dispositivo com bateria recarregável depende, por
isso, da fase de operação. Durante a descarga, a potência dissipada internamente é pequena
e a temperatura interna é ditada pela carcaça. Durante a carga, surge uma fonte interna de
calor: um carregador linear dissipa a diferença entre a tensão de entrada e a da célula
multiplicada pela corrente de carga, $P = (V_{\mathrm{in}} - V_{\mathrm{bat}})\,I$, e essa
potência se soma ao calor vindo da carcaça justamente na fase em que o limite admissível da
célula é mais baixo. É comum que o carregador monitore a temperatura da célula por um
termistor e suspenda a carga fora da faixa permitida, o que converte o limite térmico em
uma interrupção da carga, e não em dano à célula.
%% TODO: citar a folha de dados da célula LiPo (limites de 45 °C na carga e
%% 60 °C na descarga) e a do carregador BQ21040 (pino TS), usados no capítulo
%% de resultados.

%%%%%%%%%%%%%%%%%%%%%%%%%%%%%%%%%%%%%%%%%%%%%%%%%%%%%%%%%%%%%%%%%%%%%%%%%%%%%%%
\section{Síntese do capítulo}
\label{sec:fund-sintese}

Os conceitos apresentados organizam-se na mesma sequência que a informação percorre no
dispositivo. A vibração da máquina carrega assinaturas em duas escalas de frequência:
defeitos síncronos com a rotação, em harmônicas de baixa frequência, e defeitos de
rolamento, cuja informação está modulada em ressonâncias de alguns quilohertz. O caminho
mecânico até o sensor, modelado como sistema de um grau de liberdade, pode amplificar ou
distorcer essa vibração antes de ela ser medida. A amostragem e o filtro anti-rebatimento
decidem que parte do espectro chega ao processamento, e o que fica de fora não é
recuperável. A norma fixa a grandeza, a banda e os critérios com que a severidade é
avaliada, e a integração de aceleração para velocidade é a etapa em que um erro numérico
pode mudar a zona de classificação. O detector de anomalias converte descritores do sinal
em escore, e o limiar converte o escore em decisão, escolha que determina a taxa de falso
alarme; a curva ROC, a área parcial e a validação por grupos são os instrumentos que tornam
esse desempenho mensurável sem superestimá-lo. Por fim, o rádio, a bateria e o
comportamento térmico impõem os limites dentro dos quais o dispositivo pode operar.

O capítulo seguinte examina como a literatura aplicou esses conceitos a dispositivos de
baixo custo, e quais questões ela deixou em aberto.
